% RESEARCH_DRAFT_SHA256: 408ce0c4c6176653d572689a9e81e4c40b7a6eee1cb5788f0c42d754fc64f259
% Layout and source-reference wrappers expanded for the web renderer.
\documentclass[11pt]{article}
\usepackage[margin=25mm]{geometry}
\usepackage{amsthm}
\usepackage{fontspec,unicode-math}
\setmainfont{Latin Modern Roman}
\setmathfont{Latin Modern Math}
\usepackage{amsmath,amssymb,mathtools,xurl,hyperref}
\hypersetup{colorlinks=true,urlcolor=blue,linkcolor=blue}
\setcounter{secnumdepth}{0}
\setlength{\parindent}{0pt}
\setlength{\parskip}{5pt}
\setlength{\emergencystretch}{3em}
\newcommand{\sref}[2]{\hyperlink{source:#1:#2}{[S#2]}}
\newcommand{\cataloguescope}[1]{\paragraph{Source record.}#1}
\newcommand{\floor}[1]{\left\lfloor#1\right\rfloor}
\newtheorem{proposition}{Proposition}
\newtheorem{lemma}[proposition]{Lemma}
\begin{document}
\paragraph{Research manuscript.} OpenAI \texttt{gpt-6-astra}, reasoning \texttt{ultra}; 12 October 2026.
\paragraph{Preserved source body.} The complete catalogue mathematical text and sources follow with layout and source links expanded, including source qualifications and its historical preparation notes. New research begins at the section headed Research attempt.

% UnsolvedMath ID 30006187; source catalogue code OWR-14299084-010
\section{Cryptographic Equivalences for High-Energy Phenomena}\label{30006187}
{\small UnsolvedMath ID 30006187 · Computer Science}
\subsection{Definitions and mathematical statement}

\paragraph{Shared notation.}Standard number systems are $\mathbb N,\mathbb Z,\mathbb Q,\mathbb R,\mathbb C$. Specialized notation and index ranges follow the cited source and the definitions below.


A cryptographic primitive specifies a task and security experiment against an adversary. A hardness equivalence requires explicit reductions in both directions between computational problems.
\paragraph{Record scope and qualifications.}
The catalogue title is retained as metadata. Named certificates, contracts, bounds or frameworks in that title are not asserted as results of the original source extract.
The source asks for new examples and does not provide a particular phenomenon or equivalence to formulate.

Source-record limitation: Further examples of physics/cryptography equivalences are requested without a mathematical proposition.
\paragraph{Statement recorded in the supplied source.}
Find further high-energy-physics phenomena whose computational hardness is equivalent to the existence of standard cryptographic primitives, extending the known equivalence for decoding Hawking radiation. \href{https://doi.org/10.4171/owr/2025/3}{[S2]}
\paragraph{Source status.}
The supplied catalogue labels this source record open; that metadata was checked against the supplied snapshot on 2026-10-11. The label is preserved for provenance and is not an independent verification of present open status or mathematical correctness. The catalogue literature assessment is dated 2026-08-22.
\subsection{Short English statement}
Find further links between computational physics and cryptographic hardness.
\subsection{Sources}
\begin{itemize}
\item[S1] Ulam AI and the UnsolvedMath Contributors, UnsolvedMath: A Curated Collection of Open Mathematics Problems, record 30006187 (OWR-14299084-010). Catalogue curation and metadata: CC BY 4.0. The source extract is retained; definitions, English summaries and qualifications are editorial additions.
\url{https://huggingface.co/datasets/ulamai/UnsolvedMath}.
\item[S2] Cryptography (2025), Oberwolfach Reports, DOI 10.4171/owr/2025/3, pp. 141--174
\url{https://doi.org/10.4171/owr/2025/3}.
\item[S3] Brakerski, Black-Hole Radiation Decoding is Quantum Cryptography
\url{https://arxiv.org/abs/2211.05491}.
\item[S4] OWR 2025/3 Cryptography report
\url{https://ems.press/content/serial-article-files/51349}.
\item Source-specific Oberwolfach Reports licences, including CC BY-SA 4.0 where applicable, are retained; see THIRD\_PARTY\_NOTICES.md.
\end{itemize}
\paragraph{Preparation and review.}
The original source and UnsolvedMath contributors retain attribution for the source extract. Codex (GPT-6; exact serving revision not exposed) assisted with definitions, English summaries, scope notes and formatting. No new human mathematical authorship or independently verified solution is asserted. Permanent source ID 30006187 and code OWR-14299084-010 checked against the complete supplied registry and current repository. Distinct upstream records, including repeated questions and broad or damaged extracts, are retained for complete coverage of the source CS category. A repeated record is not a new independent conjecture. Complete mathematical-equivalence and current-literature review has not been performed. Preparation and final review follow contribution prompts C and D, with source gaps disclosed. Structural and rendering checks do not establish a new theorem or open status.
\label{end:30006187}

\section{Research attempt}
\subsection{Literature, source context and scope}
The full original catalogue body and source qualifications are preserved above. The relevant primary workshop passages were read in the cached publisher report. Historical claims and subsequent attributed results are distinguished from the mathematical work below. A restricted example, standard lemma or known theorem is not counted as new progress on the entire original target. Missing assumptions and repeated or rhetorical extracts are not silently turned into different solved problems.

The primary report and Brakerski's author abstract describe an existing existential equivalence involving black-hole radiation decoding and quantum cryptographic primitives, then ask for further phenomena. A newly checked 5 October 2026 preprint by Cochran and Mantri reports oracle-model query results; its abstract is context only and its full proof has not been audited here. None of these external results is claimed as new work in this attempt. No specific new physical family or desired two-way reduction is supplied by the catalogue.
\subsection{A necessary information test before a hardness equivalence}
A proposed decoding problem must distinguish computational difficulty from absence of the desired information. Let two equiprobable inputs produce accessible density matrices $\rho_0,\rho_1$ on a finite-dimensional radiation register, and write $\tau=\tfrac12\|\rho_0-\rho_1\|_1$.
\begin{proposition}
Any procedure using only that register to guess the input has average success at most $(1+\tau)/2$, irrespective of computational power. In particular, identical accessible states give success at most $1/2$.
\end{proposition}
\begin{proof}
The procedure and its final binary observation define an effect $0\preceq E\preceq I$, with probability of output zero equal to $\operatorname{Tr}(E\rho)$; any ancillary computation can be included in this measurement description. Its average success is
\[
 \tfrac12\operatorname{Tr}(E\rho_0)+\tfrac12\operatorname{Tr}((I-E)\rho_1)
 =\tfrac12+\tfrac12\operatorname{Tr}(E(\rho_0-\rho_1)).
\]
Diagonalize the Hermitian trace-zero difference as $D=D_+-D_-$ with $D_\pm\succeq0$ on orthogonal subspaces. Since $0\preceq E\preceq I$, $\operatorname{Tr}(ED)\le\operatorname{Tr}(D_+)=\|D\|_1/2=\tau$. This proves the bound. Thus success at least $1-\varepsilon$ for both inputs requires $\tau\ge1-2\varepsilon$.
\end{proof}
\subsection{A simple discarded candidate: information retained only in the inaccessible register}
Consider the isometry $V|b\rangle=|b\rangle_B|0\rangle_R$ for $b=0,1$. It preserves inner products and is implemented by putting the input in $B$ and initializing $R$ to zero. With both registers available the input is recovered by reading $B$, using a constant-size operation. Tracing out $B$ gives $\rho_0=\rho_1=|0\rangle\langle0|$ on $R$. The proposition makes decoding from $R$ impossible even without resource limits.

Consequently the inference from inaccessible information to cryptographic computational hardness fails on this candidate. A useful equivalence must impose an information-theoretic recovery promise and a computational restriction, and then prove both reductions; a generic partial trace does not provide them. This finite register example is not a model asserted to describe an actual black hole and is not a counterexample to the cited work's more specific definitions.
\section{Verification and remaining gaps}
Finite density-matrix checks verify orthogonal joint states, identical radiation reductions and the trace-distance bound. No new high-energy phenomenon, promised decoding family or cryptographic reduction has been constructed. The source is an open-ended program, so the necessary-condition analysis is auxiliary rather than original-target completion.

\section{References and source attribution}
Primary workshop report: \url{https://doi.org/10.4171/owr/2025/3}. Relevant passages read: OWR 2025/3, printed page160, Brakerski research program.

\url{https://arxiv.org/abs/2211.05491}. Read primary author abstract: existing black-hole-decoding/quantum-cryptography equivalence; full reduction proofs not audited.

\url{https://arxiv.org/abs/2610.07124}. Read5October2026 author abstract: Haar-random-oracle query bounds under stated conditions; full67page proof not audited.

All retrievals were checked on 11 October 2026. Original question attribution remains in the preserved Sources section. External theorems are attributed and not claimed as new discoveries or fully verified merely by citation.


\paragraph{Integration review.} Codex (GPT-6; exact serving revision not exposed) reviewed the written arguments and expanded layout and source-reference wrappers. Repository metadata and the production web renderer were checked using Python 3.13 and JavaScript with MathJax 3.2.2. This model-assisted review is not independent human proof verification; no PDF compilation or formal proof checking is asserted.

\section{Completion Estimate}
\noindent\textbf{COMPLETION: 0\%}
No substantiated new progress toward the entire original target is claimed. The worked results and failed route are delimited above; a known result, illustrative model, or auxiliary obstruction is not counted as original-target completion.

\section{Final status and attribution}
\textbf{Attempt status: unresolved.} This describes this research attempt, not a new ruling on the source catalogue's status or on external literature claims.
The mathematics and this self-audit were generated by OpenAI \texttt{gpt-6-astra}, reasoning \texttt{ultra}, on 12 October 2026. Codex (GPT-6; exact serving revision not exposed) orchestrated the batch, read and reviewed the written mathematical arguments, independently reran the finite checks, and contributed formatting/repository review. No human mathematical authorship, independent human proof checking, or formal verification is claimed. The historical preparation paragraphs in the reproduced catalogue body concern that body only.

Contribution prompt A was applied to the complete definition and supplied source record, followed by the adversarial checks of prompt B. Local reproduction files are \texttt{batch026/checks.py}, \texttt{computations.json}, and the per-ID audit JSON. They were not public links at preparation time. Numerical checks are finite and do not prove unrestricted results. No LaTeX engine was found locally; compilation was not performed by this research agent.

\end{document}
